\centerline{\bf \Huge Зимняя сессия по алгебре }

\begin{flushright}
        \scriptsize{\textbf{Разве ты не заметил, что способный к математике изощрен во всех науках в природе?\\
         Платон}}
        
\end{flushright}



Даны функции от \textbf{$x$}. При каких \textbf{$x$} значения функций:

а) \textit{равны нулю}

б) \textit{строго положительны}

в) \textit{не отрицательны}

г) \textit{строго отрицательны}

д) \textit{не положительны}

Баллы за каждое задание - \textbf{(а;б;в;г;д)}

\problem{\textbf{1}} $f(x)=(x+1)(x-2)+(x+1)(x-4)$

\problem{\textbf{2}} $f(x)=13|x+11|-3|x-2|$

\problem{\textbf{3}} $f(x)=\dfrac{3}{x-1}-\dfrac{-2}{6-5x}$

\problem{\textbf{4}} $f(x)=2(x-2)(5x-3)-3(9x-1)(3x-2)$

\begin{flushright} \textbf{(2;1;1;1;1)}\\ 
\end{flushright}

\problem{\textbf{5}} $f(x)=\dfrac{x-7}{3x-5}-\dfrac{2x-5}{x+4}$

\problem{\textbf{6}} $f(x)=4x^4+23x^2-35$

\problem{\textbf{7}} $f(x)=11|7-x|-5|2x+1|+4|-x+9|-2|10x+5,1|$

\begin{flushright} \textbf{(4;2;2;2;2)}\\ 
\end{flushright}




\problem{\textbf{8}} $f(x)=4|x^2-5x+6|-3|x^2+3x-28|$

\problem{\textbf{9}} $f(x)=\dfrac{1}{x-3}-\dfrac{3}{7x-\dfrac{11}{5}}+\dfrac{1}{1+2x}$

\problem{\textbf{10}} $f(x)=81(4-x)^4-16(2x-5)^4$

\begin{flushright} \textbf{(6;3;3;3;3)}\\ 
\end{flushright}

\newpage
\hspace{3mm}


\problem{\textbf{11}} $f(x)=x^4+x^3-19x^2+11x+30$

\problem{\textbf{12}} $f(x)=x^4-8x^3-31x^2-8x+1$

\begin{flushright} \textbf{(8;4;4;4;4)}\\ 
\end{flushright}


\problem{\textbf{13}} $f(x)=(x^2+2x)^2-(x+1)^2-55$

\problem{\textbf{14}} $f(x)=\dfrac{1}{x^2}+\dfrac{1}{(x+2)^2}-\dfrac{10}{9}

\begin{flushright} \textbf{(12;6;6;6;6)}\\ 
\end{flushright}


\problem{\textbf{15}} $f(x)=(2x-3)^4+20x^2-60x+31$

\problem{\textbf{16}} $f(x)=(x^2-x-1)^3+(x^2-3x+2)^3-(2x^2-4x+1)^3$

\begin{flushright} \textbf{(16;8;8;8;8)}\\ 
\end{flushright}

\centerline{\textbf{Всего 330 баллов}}








